% ============================================================
% TABLE 1B: Forecasting Studies Matrix
% ============================================================
% To change text size: modify \small below
% To change row height: modify \arraystretch{1.0} below
% X columns auto-wrap text to fit page width

\begin{landscape}
\begin{table}[htbp]
\centering
\small
\caption{Webster \& Watson Concept Matrix: Forecasting Studies}
\label{tab:forecasting-matrix}
\setlength{\tabcolsep}{3.5pt}
\renewcommand{\arraystretch}{1.0}
\begin{tabularx}{\linewidth}{@{} l X l l l X X X X X X X @{}}
\toprule
\textbf{Study} & \textbf{Design} & \textbf{Sample} & \textbf{Device} & \textbf{Loc.} & \textbf{Para.} & \textbf{Mod.} & \textbf{Pers.} & \textbf{Algorithm} & \textbf{Sens} & \textbf{Spec} & \textbf{FPR} \\
\midrule
\textbf{Vieluf et al. 2025} & Retro & n=70, 5437 d & Embrace & Wrist & Feat & M & Mix & DNN + harmonic features + diary & 82\% & 67\% & -- \\
\addlinespace
\textbf{Meisel et al. 2020} & LOSO CV & n=69, 452 szs & Empatica E4 & Wrist/Ankle & E2E & M & Glb & LSTM + 1D Conv & 51.2\% & -- & TiW 43.7\% \\
\addlinespace
\textbf{Stirling 2021} & Retro+pseudo & n=11, 136 szs & Fitbit & Wrist & Feat & M & Pat & LSTM+RF+LR (HRV feat) & -- & -- & AUC 0.74 \\
\addlinespace
\textbf{Nasseri 2021} & Retro & n=6 & Empatica E4 & Wrist & E2E & M & Tmp & LSTM 4-layer (128 hidden) & AUC 0.80 (SD 0.15) & -- & TiW 0.9--7.2h/d \\
\addlinespace
\textbf{Ode et al. 2023} & Retro & n=66, 85 szs & ECG & -- & Anom & 1 & Pat & Self-Attentive AE (attention) & 74\% & -- & 0.85/h \\
\bottomrule
\end{tabularx}
\par\medskip
\footnotesize
\textbf{Note:} ACC = accelerometer, AE = autoencoder, CNN = convolutional neural network, DNN = deep neural network, ECG = electrocardiogram, EDA = electrodermal activity, LSTM = long short-term memory, LOSO = leave-one-subject-out, PPG = photoplethysmography, Sens = sensitivity, Spec = specificity, sz(s) = seizure(s), FPR = false positive rate, TiW = time in warning, AUC = area under curve, -- = not reported. Para. = learning paradigm (Feat = feature-based, E2E = end-to-end, Anom = anomaly detection). Mod. = modality count (1 = single-modal, M = multi-modal). Pers. = personalization (Glb = global model, Pat = patient-specific, Mix = mixed, Tmp = temporal split).
\end{table}
\end{landscape}
